\section{Categorical Descent Extension}

The relational carrier in the formal framework is a functor on a schema
category. This section makes the descent condition explicit when the schema is
given by generators and path equations. The purpose is to prevent a list of
primitive relations from being mistaken for a well-defined categorical state.
Schemas as categories and graph-and-equation presentations are standard in
functorial data models \cite{spivak2012functorial,
spivakwisnesky2015relational}; TSI changes the instance codomain from sets and
functions to finite sets and binary relations.


\begin{definition}[Generator realization]
A generator realization assigns a finite nonempty set $X_A$ to every
$A\in Q_0$ and a relation $R_e\subseteq X_{s(e)}\times X_{t(e)}$ to
every $e\in Q_1$. Its free relational extension is
\[
 \widetilde F(1_A)=\Delta_{X_A},\qquad
 \widetilde F(e_n\cdots e_1)
 =R_{e_n}\circ\cdots\circ R_{e_1}.
\]
\end{definition}

\begin{proposition}[The finite-relation category]
Finite sets, binary relations, relational composition, and diagonal relations
form a category. In particular, the path formula in the preceding definition
is independent of parenthesization. This is the standard category of relations
\cite{freydscedrov1990categories}; the elementary verification is included for
completeness.
\end{proposition}

\begin{proof}
Let $R:X\nrightarrow Y$, $S:Y\nrightarrow Z$, and $T:Z\nrightarrow W$.
A pair $(x,w)$ belongs to $T\circ(S\circ R)$ exactly when there are
$y\in Y$ and $z\in Z$ with
$(x,y)\in R$, $(y,z)\in S$, and $(z,w)\in T$. Expanding
$(T\circ S)\circ R$ gives exactly the same existential condition, so
composition is associative. A pair belongs to $\Delta_Y\circ R$ exactly
when its intermediate element is forced by the diagonal to equal the target;
hence $\Delta_Y\circ R=R$. The analogous argument gives
$R\circ\Delta_X=R$. These are the category axioms.
\end{proof}

\begin{theorem}[Descent through path equations]
A generator realization descends to a functor
$F:\mathcal S\to\FinRel$ satisfying
$F([e])=R_e$ if and only if
\[
 \widetilde F(p)=\widetilde F(q)
 \qquad\text{for every }(p,q)\in\mathcal E.
\]
When it exists, the descended functor is unique.
\end{theorem}

\begin{proof}
Assume first that a descended functor $F$ exists and satisfies
$F([e])=R_e$. Let
$\pi:\operatorname{Path}(Q)\to\mathcal S$ map each path to its
equivalence class. We first verify the factorization from the generator data.
For a length-zero path, functoriality gives
$F([1_A])=\Delta_{X_A}=\widetilde F(1_A)$. If
$p=e_n\cdots e_1$ and the equality holds for paths of length $n-1$, then
functoriality and $F([e_i])=R_{e_i}$ give
\[
 F([p])=F([e_n])\circ F([e_{n-1}\cdots e_1])
 =R_{e_n}\circ\widetilde F(e_{n-1}\cdots e_1)=\widetilde F(p).
\]
Thus $\widetilde F=F\circ\pi$ by induction on path length. If
$(p,q)\in\mathcal E$, then
$\pi(p)=\pi(q)$ by construction. Therefore
$\widetilde F(p)=F(\pi(p))=F(\pi(q))=\widetilde F(q)$.

Conversely, assume equality for every declared equation. Define
$p\sim q$ exactly when $\widetilde F(p)=\widetilde F(q)$, for parallel
paths. Equality is reflexive, symmetric, and transitive. If $p\sim q$ and
$a,b$ are compatibly composable paths, functoriality of the free extension and
associativity give
\[
 \widetilde F(bpa)
 =\widetilde F(b)\circ\widetilde F(p)\circ\widetilde F(a)
 =\widetilde F(b)\circ\widetilde F(q)\circ\widetilde F(a)
 =\widetilde F(bqa).
\]
Thus $\sim$ is a path congruence containing $\mathcal E$. By minimality,
$\equiv_{\mathcal E}\subseteq\sim$. Hence
$F([p])=\widetilde F(p)$ is well-defined. The identity and composition laws
are inherited from $\widetilde F$, so $F$ is a functor and
$\widetilde F=F\circ\pi$. Every morphism of $\mathcal S$ has a path
representative, so any functor satisfying the factorization is forced to have
$F([p])=\widetilde F(p)$; uniqueness follows.
\end{proof}

\begin{proposition}[Generator criterion for natural equivalence]
Let $F,G:\mathcal S\to\FinRel$ be functors with the same object
carriers and let $\varphi_A:F(A)\to G(A)$ be bijections.
The graphs $\operatorname{Gr}(\varphi_A)$ form a natural isomorphism if and
only if, for every generator $e:A\to B$ and every $x,y$,
\[
 (x,y)\in F([e])
 \Longleftrightarrow
 (\varphi_A(x),\varphi_B(y))\in G([e]).
\]
\end{proposition}

\begin{proof}
If the family is natural, its naturality square at each generator is
$G([e])\circ\operatorname{Gr}(\varphi_A)
=\operatorname{Gr}(\varphi_B)\circ F([e])$. Expanding graph
composition gives the displayed biconditional: one direction follows by
transporting a related pair, and the other uses surjectivity of
$\varphi_B$.

Conversely, the biconditional transports each generator relation. It also
holds for identities because a bijection transports a diagonal relation to a
diagonal relation. For composable paths, membership in a relational composite
is witnessed by an intermediate entity; applying the criterion to both factors
and using bijectivity transports that witness. Induction on path length gives
the criterion for every path. Since both functors satisfy the same path
equations, it is independent of the chosen representative. The corresponding
naturality equality therefore holds for every schema morphism. Each graph is
an isomorphism in $\FinRel$, so the natural transformation is a natural
isomorphism.
\end{proof}

If a declared path equation is not satisfied by the generator relations, the
assignment remains a valid functor on the free path category but is not a
state over the presented schema. Thus categorical validity is an additional
mathematical gate, not an interpretive label attached after prediction.
The extension does not assert that a learner discovers the correct schema or
that generator-level agreement is sufficient without checking the equations.
